\chapter{Related Work}
\label{sec:related_work}

\section{Organising the Literature}
A related-work chapter should help the reader understand why the study makes
particular methodological choices. One possible organisation separates
retrieval methods, evaluation resources and analysis procedures. Another
follows the research questions directly. The choice should make comparisons
clear rather than simply reproduce the order in which papers were read.
For each relevant study, the discussion should identify its problem, evidence,
and relationship to the present work.

The citations in this demonstration illustrate formatting only. The laboratory
website entry~\cite{RSL} demonstrates a web address in a bibliography. The journal
entry~\cite{exampleArticle} and conference entry~\cite{exampleConference} are
fictional and are not sources for the methodological discussion. In a real
thesis, replace these entries with verified publications and ensure that each
citation supports the sentence in which it appears.

\section{Comparing Experimental Settings}
Two studies may use similar ranking methods while evaluating substantially
different tasks. Differences in collection size, query selection, judgement
coverage or metric definition can affect the interpretation of their scores.
A careful comparison therefore reports the experimental setting alongside the
method. It should avoid treating values from unrelated collections as if they
were measurements on the same scale under identical conditions.

A compact comparison table can summarise these settings, but the surrounding
text still needs to explain their significance. For instance, a collection
containing short factual queries may not reveal the same behaviour as a
collection of exploratory requests. Such differences can motivate an
additional experiment without implying that either earlier study was flawed.

% Demonstration-only break; remove for natural pagination in a real thesis.
\newpage
\section{Identifying a Research Gap}
A research gap should be stated in terms that can be investigated. Saying that
a topic has received limited attention is less specific than identifying an
unresolved comparison or an untested assumption. In this example, the proposed
gap concerns whether a change in average retrieval quality is accompanied by
consistent behaviour across query groups. That statement leads to observable
quantities and a concrete analysis plan.

The scope of the claim should match the literature that was examined. A thesis
can explain its search strategy, inclusion criteria and stopping point without
claiming that every relevant publication has been found. It can also separate
a practical implementation gap from a conceptual limitation. These distinctions
help prevent an overly broad novelty claim from obscuring a useful, narrower
contribution.

\section{Position of the Present Example}
The proposed comparison uses an established baseline as a reference point and
changes a single component. The role of related work is to justify that baseline,
explain reasonable alternatives and identify possible confounding factors.
A stronger or newer method is not automatically the most informative baseline
if its data preparation or computational requirements cannot be matched.

The next chapter translates these considerations into a procedure. It defines
the inputs, the two configurations and the aggregation of per-query scores.
This transition connects the literature review with the actual experiment,
so that the reader can see how earlier observations influenced the study
design rather than encountering an unrelated list of publications.
